\documentclass[14pt]{article}
\usepackage{amsmath, amssymb, amsfonts}
\usepackage{enumerate}
\usepackage{hyperref}

\parindent 0px
%beamer: presentation
%exams: quizes

\begin{document}

	\section{Hello \LaTeX Math}
	Hello this is  not  my first \LaTeX document. % a hard return.
	
	A rectangle of sides $(x + 1)$ and $(x + 3)$\\ %soft return with \\
	The equations that is sol long long polinomial and quadratic equation $A(x) = x^2 + 6x + 3$ gives the area of the rectangle.\\
	The equations that is sol long long polinomial and quadratic equation ${A(x) = x^2 + 6x + 3}$ gives the area of the rectangle.\\
	The equations that is sol long long polinomial and quadratic equation $$A(x) = x^2 + 6x + 3$$ gives the area of the rectangle.
	
	\section{$$^{\text{sup}}_{\text{sub}}$$scripts}
	
	$$2_{1_{2_{3_4}}}$$	
	
	$$3^\textbf{HELLO superscript}$$
	
	\section{Trig functions}
	
	\[\sin x\]
	$$\tan45^0=1$$
	$$\arcsin\beta$$
	$$\cos(2n\pi t/T)$$
	
	\section{Logarithmic}
	
	$$\ln e$$
	$$\log100 = 2$$
	$$\log_a a^k = k$$
	
	\section{Roots}
	
	$$\sqrt[x]{y}$$
	$$\sqrt{a^2+b^2}=c$$
	
	\section{Fractions}
	$$x_1,x_2=\frac{-b+\sqrt{b^2-4ac}}{2a}$$
	
	About $\frac12$ of the glass is full	
		
	about $\displaystyle\frac12$ of the glass is empty
			
	about $\dfrac{1}{2}$ of the glass is not used
	
	about $\dfrac{1}{2}$ of the glass is used
	
	\section{Bracketting}
	
	The distributive property states that $a(b+c) = ab + ac$ for all $a, b, c \in \mathbb{R}$.
	
	The equivalent class is $a$ is $[a]$.\\
	The set $A$ is defined to be $\{1,2,3,4\}$
	$$f(x)=a_0 + \sum^{\infty}_{n=1} \left(
		a_n\cos\left(\frac{2\pi nt}{T}\right) + 
		b_n\sin\left(\frac{2\pi nt}{T}\right) 
	\right)$$
	
	Conversion rate as of Sept, 2026: $\text{UGX}4000\approx\$1\text{(USD)}$
	
	$$b_n=\frac{2}{T}\int_Tf(x)\sin\frac{2\pi nt}{T}$$
	$$a_n=\frac{2}{T}\int_Tf(x)\cos\left[\frac{2\pi nt}{T}\right]$$
	$$b_n=\frac{2}{T}\int_Tf(x)\sin\left\langle\frac{2\pi nt}{T}\right\rangle$$
	$$\left.\frac{dx}{dy}\right|_{x=5}$$
	
	\section{Tabulation}
	
	\begin{tabular}{|lc||c|}
	\hline
	A & B & X \\
	\hline
	0 & 0 & 1 \\
	0 & 1 & 1 \\
	1 & 0 & 1 \\
	1 & 1 & 0 \\
	\hline
	
	\end{tabular}
	
	\begin{table}
		\centering
		\def \arraystretch{1.5}
	
		\begin{tabular}{|c|c|c|}
		\hline
			A & B & C \\
		\hline
			1 & $\frac{1}{5}$ & 5\\ \hline
			2 & 4 & 12\\ \hline
			12 & 12 & $\dfrac{1+\dfrac{2}{3x^2}}{5n!-2}$\\
		\hline
		\end{tabular}
		\caption{A random tabular data}
	\end{table}
	
	\begin{table}
		\centering		
		
		\caption{A relationship between parameter, argument and function}
		\def \arraystretch{1.25}
		
		\begin{tabular}{|l|p{5cm}|}
		\hline
		var, fn & Description \\ 
		\hline
		$x,f(x)$ & $f(x)$ is the reusable piece of math/code that accepts uses the parameter $x$. A value that is pass to the function eg in $f(5), 5$ is called the \textbf{argument}.\\
		\hline
		
		\end{tabular}
	\end{table}
	
	\section{Arrays}
	
	\begin{align}
		x^2 + 2x -5 &= 0\quad \text{is a quadratic equation}\\
		x+5y&=4
	\end{align}
	
	
	\begin{align*}
		x^2 + 2x -5 &= 0\quad \text{is a quadratic equation}\\
		x+5y&=4
	\end{align*}
	
	
	\begin{align}
		x^2 + 2x -5 &= 0\quad \text{is a quadratic equation}\\
		x+5y&=4
	\end{align}
	
	\section{Itemization}
	
	\begin{enumerate}
		\item Ruler
		\item Pen
		\item Exam permit
		\item Calculator
		\item etc.
	\end{enumerate}
	
	
	\begin{itemize}
		\item Ruler
		\item Pen
			\begin{itemize}
				\item[$\times$] Biro pen
				\item[$\checkmark$] Fountain pen
				\item[$\times$] Coloured pen
			\end{itemize}
		\item Exam permit
		\item Calculator
		\item etc.
	\end{itemize}
	
	\begin{enumerate}[A:-]\setcounter{enumii}{5}
		\item Ruler
		\item Pen
			\begin{enumerate}
				\item Biro pen
				\item Fountain pen
				\item Coloured pen
			\end{enumerate}
		\item Exam permit
		\item Calculator
		\item etc.
	\end{enumerate}
	
	\section{Text formating}
	
	This sentence contains \textit{italicized}, \textbf{bold}, \textsc{small caps}, \textsf{serif}, \texttt{monospace}, \textsl{slanted/oblique} text.
	
	View my Github page on \texttt{https://github.com/waraton}.
	
	Acess this course video at \url{https://youtube.com?q=latex-course-by-Michelle-Krummel-freecodecamp}.	
	
	Acess this course video at \href{https://youtube.com?q=latex-course-by-Michelle-Krummel-freecodecamp}{FreeCodeCamp on YouTube}.
	
	\subsection{sizes}
	\begin{small}\end{small}\\
	\begin{smaller}\end{smaller}\\
	\begin{tiny}\end{tiny}\\
	\begin{big}\end{small}\\
	\begin{small}\end{small}\\
	\begin{small}\end{small}\\
	\begin{small}\end{small}\\
	\begin{small}\end{small}\\
	\begin{small}\end{small}\\
	\begin{small}\end{small}\\
	\begin{small}\end{small}\\
	\begin{small}\end{small}\\
	\begin{small}\end{small}\\
	\begin{small}\end{small}\\
	\begin{small}\end{small}\\
	\begin{small}\end{small}\\
	\begin{small}\end{small}\\
		
\end{document}
